\chapter{Couple Bubble: Gelembung Pasangan yang Melindungi}

\begin{insightbox}
\textit{``Couple Bubble bukanlah tentang isolasi dari dunia, melainkan tentang menciptakan zona aman di mana kalian berdua bisa saling melindungi. Ini adalah pakta mutual protection: \highlight{aku menjagamu, kamu menjagaku, dan kita menjaga \textit{kita} bersama}.''} --- Dr. Stan Tatkin
\end{insightbox}

\section{Memahami Konsep Couple Bubble}

\subsection{Apa itu Couple Bubble?}

\keyterm{Couple Bubble} adalah sebuah konsep yang dikembangkan oleh Dr. Stan Tatkin untuk menggambarkan \highlight{ruang psikologis dan emosional yang diciptakan oleh pasangan} untuk saling melindungi. Bayangkan sebuah gelembung transparan yang membungkus kalian berdua --- di dalamnya, kalian merasa aman, diterima, dan saling mendukung.

\begin{praktikbox}[Definisi Couple Bubble]
Couple Bubble adalah \textbf{kesepakatan implisit dan eksplisit} antara dua orang untuk:
\begin{itemize}
    \item Menempatkan keamanan hubungan sebagai prioritas utama
    \item Menjadi tempat perlindungan satu sama lain dari ancaman dunia luar
    \item Berkomitmen untuk saling menjaga, meski dalam konflik
    \item Membangun sistem kekebalan bersama terhadap tekanan eksternal
\end{itemize}
\end{praktikbox}

\subsection{Mengapa Couple Bubble Penting?}

Secara evolusioner, manusia adalah makhluk \textit{pair-bonding} --- kita butuh ikatan dengan pasangan untuk bertahan hidup dan berkembang biak. Otak kita \textit{wired} (terkabel) untuk:

\begin{enumerate}
    \item \textbf{Mencari keamanan} dalam hubungan dekat
    \item \textbf{Menggunakan pasangan} sebagai ``secure base'' untuk menjelajahi dunia
    \item \textbf{Saling menenangkan} saat stres (co-regulation)
\end{enumerate}

\begin{center}
\begin{tikzpicture}[
    bubble/.style={draw=primary, line width=3pt, rounded corners=20pt, 
                   fill=lightbg, minimum width=12cm, minimum height=7cm, align=center}
]
    % Outer circle representing the bubble
    \draw[primary, line width=4pt, dashed, fill=primary!5] (0,0) circle (5.5cm);
    
    % Couple inside
    \node[circle, draw=sagegreen, line width=2pt, fill=sagegreen!20, minimum size=2.5cm] (p1) at (-2,0) {};
    \node[font=\large\bfseries, sagegreen] at (p1) {AKU};
    
    \node[circle, draw=secondary, line width=2pt, fill=secondary!20, minimum size=2.5cm] (p2) at (2,0) {};
    \node[font=\large\bfseries, secondary] at (p2) {KAMU};
    
    % Connection
    \draw[<->, line width=3pt, primary] (p1) -- (p2) node[midway, above, font=\small\bfseries] {PROTEKSI MUTUAL};
    
    % External threats
    \node[circle, draw=accent, line width=1pt, fill=accent!10, minimum size=1.2cm] (t1) at (-4.5,3) {\small Kerja};
    \node[circle, draw=accent, line width=1pt, fill=accent!10, minimum size=1.2cm] (t2) at (4.5,3) {\small Mertua};
    \node[circle, draw=accent, line width=1pt, fill=accent!10, minimum size=1.2cm] (t3) at (-4.5,-3) {\small Teman};
    \node[circle, draw=accent, line width=1pt, fill=accent!10, minimum size=1.2cm] (t4) at (4.5,-3) {\small HP};
    
    % Arrows showing protection
    \draw[->, line width=2pt, accent, dashed] (t1) -- (-3.5,2);
    \draw[->, line width=2pt, accent, dashed] (t2) -- (3.5,2);
    \draw[->, line width=2pt, accent, dashed] (t3) -- (-3.5,-2);
    \draw[->, line width=2pt, accent, dashed] (t4) -- (3.5,-2);
    
    % Bubble label
    \node[font=\Large\bfseries, primary] at (0,4.5) {COUPLE BUBBLE};
    \node[font=\small, warmgray] at (0,-4.5) {\textit{Zona Aman Bersama}};
\end{tikzpicture}
\end{center}

\subsection{Perbedaan Couple Bubble dan Ketergantungan Tidak Sehat}

\begin{warningbox}[Couple Bubble $\neq$ Codependency]
\textbf{Couple Bubble (Sehat):}
\begin{itemize}
    \item Dua individu utuh yang memilih untuk saling melindungi
    \item Batasan dengan dunia luar jelas tapi fleksibel
    \item Kedua pihak tetap bisa berfungsi secara independen
    \item Proteksi bersifat saling, bukan satu arah
\end{itemize}

\textbf{Codependency (Tidak Sehat):}
\begin{itemize}
    \item Satu pihak ``menyelamatkan'' pihak lain secara terus-menerus
    \item Batasan kabur atau tidak ada
    \item Kehilangan identitas individual
    \item Dinamasa yang tidak setara (caregiver-pasien)
\end{itemize}
\end{warningbox}

\section{Tiga Prinsip Fundamental Couple Bubble}

\subsection{Prinsip 1: Menempatkan Hubungan di Atas Segalanya}

Ini adalah pilar paling fundamental. Artinya: \highlight{kepentingan hubungan kalian berdua lebih penting daripada kebutuhan individu, pekerjaan, keluarga besar, atau teman}.

\begin{tipbox}[Contoh Prinsip 1 dalam Praktik]
\textbf{Situasi:} Pasanganmu memiliki hari yang buruk dan butuh bicara, tapi kamu sedang sibuk dengan deadline kerja.

\textbf{Tanpa Couple Bubble:} ``Aku sibuk, kita bicara nanti saja. Kerjaanku lebih penting.''

\textbf{Dengan Couple Bubble:} ``Aku melihat kamu butuhku sekarang. Beri aku 10 menit untuk menyelesaikan ini, lalu aku akan matikan laptop dan kita bicara. Hubungan kita lebih penting daripada deadline ini.''
\end{tipbox}

\subsection{Prinsip 2: Menjadi \textit{Go-To Person} Satu Sama Lain}

\keyterm{Go-to person} adalah orang pertama yang kamu hubungi saat:
\begin{itemize}
    \item Ada kabar baik (\textit{capitalization})
    \item Ada masalah atau stres
    \item Butuh pendapat atau saran
    \item Merasa takut, sedih, atau bingung
\end{itemize}

\begin{praktikbox}[Mengapa Go-To Person Penting?]
Otak kita membutuhkan \textbf{secure base} --- tempat yang aman untuk ``pulang'' secara emosional. Ketika pasanganmu menjadi go-to person-mu:
\begin{enumerate}
    \item Sistem sarafmu menjadi lebih rileks
    \item Kamu merasa diterima sepenuhnya
    \item Kepercayaan diri meningkat karena ada ``safety net''
    \item Resiliensi terhadap stres eksternal meningkat
\end{enumerate}

\textit{Couple yang kuat adalah couple di mana masing-masing berkata: ``Aku tidak perlu orang lain, aku punya dia/dia.''}
\end{praktikbox}

\subsection{Prinsip 3: Melindungi Satu Sama Lain dari Ancaman}

Ancaman bisa datang dari dua arah:

\begin{center}
\renewcommand{\arraystretch}{1.6}
\begin{tabular}{|>{\raggedright}p{3cm}|>{\raggedright}p{5.5cm}|>{\raggedright\arraybackslash}p{5.5cm}|}
\hline
\rowcolor{primary!20}
\textbf{Jenis Ancaman} & \textbf{Contoh} & \textbf{Cara Melindungi} \\
\hline
\textbf{Eksternal} & Tekanan kerja, kritik mertua, gosip teman, distraksi teknologi & Bersatu menghadapi, membuat batasan, saling menenangkan \\
\hline
\textbf{Internal} & Kebiasaan buruk sendiri, insecure attachment, trauma masa lalu & Mengingatkan dengan lembut, memberi reassurance, tidak menyalahkan \\
\hline
\textbf{Interpersonal} & Konflik antar pasangan, misunderstanding, resentmen & Repair cepat, jangan biarkan mengendap, prioritaskan hubungan \\
\hline
\end{tabular}
\end{center}

\section{Self-Assessment: Seberapa Kuat Couple Bubble-mu?}

\begin{exercisebox}[Kuesioner Kekuatan Couple Bubble]
\textbf{Petunjuk:} Jawablah dengan jujur menggunakan skala berikut:
\begin{center}
\textbf{1} = Sangat tidak setuju \quad | \quad \textbf{2} = Tidak setuju \quad | \quad \textbf{3} = Netral \quad | \quad \textbf{4} = Setuju \quad | \quad \textbf{5} = Sangat setuju
\end{center}

\subsection*{Dimensi 1: Prioritas Hubungan}

\begin{enumerate}
    \item Saat ada konflik antara kebutuhanku dan kebutuhan hubungan, aku biasanya memprioritaskan hubungan.
    \item Aku sering mengorbankan waktu pribadi demi menghabiskan waktu bersama pasangan.
    \item Saat pasangan butuhku, aku bisa meletakkan pekerjaan atau hobi sejenak.
    \item Aku merasa hubungan kami adalah aset paling berharga dalam hidupku.
\end{enumerate}

\subsection*{Dimensi 2: Saling Melindungi}

\begin{enumerate}
    \setcounter{enumi}{4}
    \item Aku merasa aman membagikan kelemahanku pada pasangan tanpa takut dihakimi.
    \item Pasangan sering membelaku saat aku dikritik orang lain.
    \item Kami berdua saling mendukung saat salah satu stres dengan hal luar.
    \item Aku percaya pasangan akan selalu ``ada untukku'' saat aku benar-benar butuh.
\end{enumerate}

\subsection*{Dimensi 3: Batasan dengan Dunia Luar}

\begin{enumerate}
    \setcounter{enumi}{8}
    \item Kami punya aturan bersama tentang penggunaan ponsel saat bersama.
    \item Kami bisa menolak permintaan keluarga/teman yang mengganggu waktu berdua.
    \item Aku tidak membicarakan masalah hubungan dengan orang lain tanpa persetujuan pasangan.
    \item Kami menjaga privasi hubungan dari interferensi orang luar.
\end{enumerate}

\subsection*{Dimensi 4: Resiliensi dalam Konflik}

\begin{enumerate}
    \setcounter{enumi}{12}
    \item Saat bertengkar, kami tetap merasa ``kita dalam tim yang sama.''
    \item Kami cepat melakukan repair setelah konflik.
    \item Bahkan saat marah, aku tidak mengancam akan pergi atau mengakhiri hubungan.
    \item Aku merasa kami bisa melewati apa pun asal bersama-sama.
\end{enumerate}

\subsection*{Skoring dan Interpretasi}

\begin{center}
\renewcommand{\arraystretch}{1.4}
\begin{tabular}{|c|c|l|}
\hline
\rowcolor{sagegreen!20}
\textbf{Total Skor} & \textbf{Kategori} & \textbf{Interpretasi} \\
\hline
56--70 & \textbf{Sangat Kuat} & Couple Bubble kalian sangat solid. Pertahankan dengan ritual harian. \\
\hline
42--55 & \textbf{Kuat} & Ada fondasi yang baik. Beberapa area bisa ditingkatkan. \\
\hline
28--41 & \textbf{Perlu Perhatian} & Bubble mulak berlubang. Segera lakukan perbaikan bersama. \\
\hline
14--27 & \textbf{Lemah} & Bubble dalam bahaya. Intervensi intensif diperlukan. \\
\hline
\end{tabular}
\end{center}
\end{exercisebox}

\section{Ancaman terhadap Couple Bubble}

\subsection{Ancaman Eksternal}

\begin{warningbox}[Musuh dari Luar]
\begin{enumerate}
    \item \textbf{Pekerjaan yang Menyerbu}
    \begin{itemize}
        \item Email dan chat kerja yang di-check terus-menerus
        \item Membawa pulang stres dan frustrasi kantor
        \item Mengutamakan deadline daripada pasangan
    \end{itemize}
    
    \item \textbf{Keluarga Besar (Mertua)}
    \begin{itemize}
        \item Interferensi dalam keputusan pasangan
        \item Kritik terhadap pasanganmu
        \item Menempatkan loyalty keluarga di atas pasangan
    \end{itemize}
    
    \item \textbf{Teman dan Sosial Media}
    \begin{itemize}
        \item Menghabiskan lebih banyak waktu dengan teman daripada pasangan
        \item Membagi masalah hubungan di media sosial
        \item Teman yang meremehkan pasanganmu
    \end{itemize}
    
    \item \textbf{Smartphone dan Digital Distraction}
    \begin{itemize}
        \item Phubbing (phone + snubbing) --- mengabaikan pasangan karena ponsel
        \item Mengobrol dengan orang lain saat waktu bersama pasangan
        \item Tidak ada ``sacred space'' bebas gadget
    \end{itemize}
\end{enumerate}
\end{warningbox}

\subsection{Ancaman Internal}

\begin{center}
\begin{tikzpicture}[
    threat/.style={draw=#1, line width=2pt, rounded corners=10pt,
                   fill=#1!10, minimum width=5.5cm, minimum height=3cm, align=left, text width=5cm}
]
    % Internal threats
    \node[threat=accent] (i1) at (-4,2) {
        \textbf{\textcolor{accent}{Resentment}}\\[0.1cm]
        \small Menyimpan amarah\\
        \small yang tidak diungkapkan\\
        \small Memendam kekesalan
    };
    
    \node[threat=accent] (i2) at (4,2) {
        \textbf{\textcolor{accent}{Rahasia}}\\[0.1cm]
        \small Menyembunyikan\\
        \small hal penting\\
        \small Keuangan, komunikasi\\
        \small dengan orang lain
    };
    
    \node[threat=accent] (i3) at (-4,-2) {
        \textbf{\textcolor{accent}{Attachment Issues}}\\[0.1cm]
        \small Avoidance (Island)\\
        \small Anxiety (Wave)\\
        \small Membuat pasangan\\
        \small merasa tidak aman
    };
    
    \node[threat=accent] (i4) at (4,-2) {
        \textbf{\textcolor{accent}{Prioritas Salah}}\\[0.1cm]
        \small Anak di atas pasangan\\
        \small Pekerjaan nomor satu\\
        \small Hobi menguasai waktu
    };
    
    % Center bubble
    \node[circle, draw=primary, line width=3pt, fill=lightbg, minimum size=2.5cm, align=center, font=\bfseries\color{primary}] (center) at (0,0) {BUBBLE\\TERANCAM};
    
    % Arrows
    \draw[->, line width=1.5pt, accent] (i1) -- (center);
    \draw[->, line width=1.5pt, accent] (i2) -- (center);
    \draw[->, line width=1.5pt, accent] (i3) -- (center);
    \draw[->, line width=1.5pt, accent] (i4) -- (center);
\end{tikzpicture}
\end{center}

\begin{praktikbox}[Mengenali Tanda-Tanda Bubble Bermasalah]
\textbf{Sinyal Awal (Kuning):}
\begin{itemize}
    \item JarangQuality time tanpa gangguan
    \item Mengobrol dengan orang lain tentang pasangan (komplain)
    \item Merasa lebih nyaman sendiri atau dengan teman
\end{itemize}

\textbf{Sinyal Bahaya (Merah):}
\begin{itemize}
    \item Berpikir ``Aku lebih bahagia tanpa dia/dia''
    \item Menyimpan rahasia besar dari pasangan
    \item Fantasi tentang orang lain atau hidup tanpa pasangan
    \item Tidak merasa ``tim'' saat menghadapi masalah
\end{itemize}
\end{praktikbox}

\section{Membangun Couple Bubble: Panduan Langkah demi Langkah}

\subsection{Langkah 1: Komitmen Bersama}

Buat kesepakatan formal bersama. Tuliskan dan tanda tangani.

\begin{exercisebox}[Template Komitmen Couple Bubble]
\textbf{KOMITMEN COUPLE BUBBLE}

Kami, \underline{\hspace{4cm}} dan \underline{\hspace{4cm}}, 
sepakat untuk membangun dan menjaga Couple Bubble kami dengan komitmen sebagai berikut:

\textbf{1. PRIORITAS HUBUNGAN}
\begin{itemize}
    \item[$\square$] Hubungan kami adalah prioritas utama setelah kesehatan dan keselamatan
    \item[$\square$] Kami akan selalu memikirkan ``apa yang terbaik untuk kita'' sebelum mengambil keputusan besar
    \item[$\square$] Kami berkomitmen untuk menghabiskan waktu berkualitas minimal \underline{\hspace{2cm}} per minggu
\end{itemize}

\textbf{2. PROTEKSI MUTUAL}
\begin{itemize}
    \item[$\square$] Aku akan selalu ada untukmu saat kamu butuh, meski aku sedang sibuk
    \item[$\square$] Kami akan saling melindungi dari interferensi eksternal (kerja, keluarga, teman)
    \item[$\square$] Kami tidak akan membicarakan masalah kami dengan orang luar tanpa persetujuan bersama
\end{itemize}

\textbf{3. KOMUNIKASI}
\begin{itemize}
    \item[$\square$] Kami akan berbicara dengan jujur tentang kebutuhan dan kekhawatiran
    \item[$\square$] Kami akan melakukan repair secepatnya setelah konflik
    \item[$\square$] Kami akan menggunakan ``stress-reducing conversation'' minimal 20 menit per hari
\end{itemize}

\textbf{4. BATASAN}
\begin{itemize}
    \item[$\square$] Kami akan punya ``no-phone zone'' dan ``no-phone time'' bersama
    \item[$\square$] Kami bisa menolak permintaan yang mengganggu waktu berdua
    \item[$\square$] Privasi hubungan kami dijaga dari media sosial
\end{itemize}

\vspace{0.5cm}
Tanda tangan:\\[0.3cm]
\underline{\hspace{6cm}} \quad Tanggal: \underline{\hspace{3cm}}\\[0.2cm]
(\textit{Pasangan 1})\\[0.5cm]
\underline{\hspace{6cm}} \quad Tanggal: \underline{\hspace{3cm}}\\[0.2cm]
(\textit{Pasangan 2})
\end{exercisebox}

\subsection{Langkah 2: Membangun Ritual Koneksi}

\begin{tipbox}[Ritual Harian untuk Bubble Kuat]
\textbf{Pagi (5-10 menit):}
\begin{itemize}
    \item Sapaan hangat dengan kontak mata dan pelukan
    \item Tanya: ``Ada yang perlu kuketahui untuk hari ini?''
    \item Janji akan kembali dengan selamat
\end{itemize}

\textbf{Siang (1-2 menit):}
\begin{itemize}
    \item Pesan singkat: ``Thinking of you'' atau ``Semangat!''
    \item Tidak untuk diskusi panjang, hanya penanda ``aku di sini''
\end{itemize}

\textbf{Malam (20+ menit):}
\begin{itemize}
    \item Stress-reducing conversation (dengarkan tanpa solusi)
    \item Physical affection: pegang tangan, peluk, ciuman
    \item Ritual tidur: bersamaan, sentuhan sebelum tidur
\end{itemize}
\end{tipbox}

\subsection{Langkah 3: Membuat Batasan yang Jelas}

\begin{center}
\renewcommand{\arraystretch}{1.5}
\begin{tabular}{|>{\raggedright}p{4cm}|>{\raggedright}p{5cm}|>{\raggedright\arraybackslash}p{5cm}|}
\hline
\rowcolor{sagegreen!20}
\textbf{Area} & \textbf{Batasan yang Sehat} & \textbf{Implementasi} \\
\hline
Penggunaan ponsel & Tidak ponsel saat makan dan 1 jam sebelum tidur & Taruh di tempat tertentu, gunakan mode DND \\
\hline
Waktu kerja & Kerja tetap di kantor/ruang kerja, tidak ``bawa pulang'' & Buat ritual transisi: mandi, ganti baju sebelum quality time \\
\hline
Keluarga besar & Keputusan penting dibahas berdua dulu & Komunikasikan pada keluarga: ``Kami akan diskusi dulu'' \\
\hline
Teman & Waktu dengan teman tidak menggantikan date night & Jadwalkan teman di luar waktu prioritas pasangan \\
\hline
Media sosial & Tidak posting masalah hubungan & Agreement: tanya izin sebelum posting foto pasangan \\
\hline
\end{tabular}
\end{center}

\section{Maintenance Harian untuk Couple Bubble}

\begin{praktikbox}[Checklist Pemeliharaan Bubble]
\textbf{Setiap Hari (minimal):}
\begin{itemize}
    \item[$\square$] Morning greeting dengan kontak fisik (minimal 6 detik pelukan)
    \item[$\square$] 20 menit stress-reducing conversation
    \item[$\square$] Satu momen appreciation: ``Aku suka karena kamu...''
    \item[$\square$] Physical touch sebelum tidur
    \item[$\square$] Check-in singkat via pesan saat terpisah
\end{itemize}

\textbf{Setiap Minggu:}
\begin{itemize}
    \item[$\square$] Date night 2+ jam tanpa gangguan (no phone!)
    \item[$\square$] State of the Union meeting: ``Gimana hubungan kita minggu ini?''
    \item[$\square$] Satu aktivitas bersama yang menyenangkan
    \item[$\square$] Review: Apa yang berhasil? Apa yang perlu diperbaiki?
\end{itemize}

\textbf{Setiap Bulan:}
\begin{itemize}
    \item[$\square$] Review dan update Couple Agreement
    \item[$\square$] Planning untuk bulan depan (date, goals)
    \item[$\square$] Revisit Love Maps (update pengetahuan tentang pasangan)
\end{itemize}
\end{praktikbox}

\section{Perbaikan Darurat Saat Bubble Bocor}

\begin{warningbox}[Tanda Bubble dalam Bahaya]
\begin{enumerate}
    \item Salah satu pihak mengatakan ``Aku sudah tidak yakin dengan kita''
    \item Terjadi pengkhianatan kepercayaan (kebohongan besar, selingkuh emosional/fisik)
    \item Resentmen yang mendalam dan mengendap
    \item Kehilangan ikatan intem (tidak merasa ``kita'' lagi)
    \item Ancaman perceraian/pengakhiran hubungan sering muncul
\end{enumerate}
\end{warningbox}

\subsection{Protokol Perbaikan Darurat}

\begin{center}
\begin{tikzpicture}[
    stepbox/.style={draw=primary, line width=2pt, rounded corners=10pt,
                 fill=lightbg, minimum width=13cm, minimum height=2.5cm, align=left, text width=12cm}
]
    \node[stepbox] (s1) at (0,6) {
        \textbf{\large LANGKAH 1: STOP}\\[0.1cm]
        \textbullet\ Jangan lanjutkan pola yang merusak\\
        \textbullet\ Akui bahwa Bubble bermasalah\\
        \textbullet\ Sepakat untuk ``mengalahkan musuh bersama'' bukan satu sama lain
    };
    
    \node[stepbox] (s2) at (0,3) {
        \textbf{\large LANGKAH 2: ASSESS}\\[0.1cm]
        \textbullet\ Diskusikan: Apa yang membuat Bubble kita bocor?\\
        \textbullet\ Identifikasi ancaman eksternal dan internal\\
        \textbullet\ Gunakan kuesioner di bab ini untuk objektivitas
    };
    
    \node[stepbox] (s3) at (0,0) {
        \textbf{\large LANGKAH 3: REPAIR}\\[0.1cm]
        \textbullet\ Minta maaf untuk kontribusi masing-masing terhadap keretakan\\
        \textbullet\ Buat rencana konkret perbaikan dengan timeline\\
        \textbullet\ Jika perlu, cari bantuan profesional (counselor)
    };
    
    \node[stepbox] (s4) at (0,-3) {
        \textbf{\large LANGKAH 4: REBUILD}\\[0.1cm]
        \textbullet\ Kembali ke ritual dasar (morning/evening ritual)\\
        \textbullet\ Tingkatkan quality time secara bertahap\\
        \textbullet\ Rayakan small wins dalam perbaikan
    };
    
    % Arrows
    \draw[->, line width=3pt, sagegreen] (s1) -- (s2);
    \draw[->, line width=3pt, sagegreen] (s2) -- (s3);
    \draw[->, line width=3pt, sagegreen] (s3) -- (s4);
\end{tikzpicture}
\end{center}

\begin{tipbox}[Teknik Repair Cepat]
\textbf{Jika baru terjadi konflik atau kesalahpahaman:}

\begin{enumerate}
    \item \textbf{The Pause} --- Beri jeda 20-30 menit agar sistem saraf tenang
    \item \textbf{The Approach} --- Salah satu mendekati dengan sikap terbuka
    \item \textbf{The Validation} --- ``Aku mengerti kamu merasa...'' (tanpa pembelaan)
    \item \textbf{The Ownership} --- ``Maaf aku sudah...'' (akui bagianmu)
    \item \textbf{The Reconnect} --- Sentuhan fisik (pegang tangan/peluk) untuk mengaktifkan oxytocin
    \item \textbf{The Plan} --- Diskusikan apa yang bisa dilakukan berbeda
\end{enumerate}

\textit{Ingat: Tujuan bukan memenangkan argumen, tapi memenangkan hubungan.}
\end{tipbox}

\section{Integrasi dengan Gaya Koneksi (Attachment Styles)}

Setiap tipe attachment membutuhkan pendekatan berbeda dalam membangun Bubble:

\subsection{Anchor (Secure) Membangun Bubble}

\begin{praktikbox}[Kekuatan Anchor dalam Couple Bubble]
\textbf{Yang Natural:}
\begin{itemize}
    \item Mudah membuat pasangan merasa prioritas
    \item Balance antara dekat dan memberi ruang
    \item Komunikasi terbuka tentang kebutuhan
\end{itemize}

\textbf{Yang Perlu Diperhatikan:}
\begin{itemize}
    \item Jangan jadi ``rescuer'' yang selalu mengorbankan diri
    \item Pastikan pasangan insecure juga berinvestasi
    \item Tetap jaga ritual meski merasa ``aman''
\end{itemize}
\end{praktikbox}

\subsection{Island (Avoidant) Membangun Bubble}

\begin{warningbox}[Tantangan Island]
\textbf{Hambatan internal Island:}
\begin{itemize}
    \item ``Couple Bubble terasa seperti jebakan/kurungan''
    \item ``Aku butuh ruang, jangan terlalu dekat''
    \item Sulit meminta bantuan atau menunjukkan kebutuhan
    \item Tendensi untuk mengutamakan pekerjaan/hobi di atas hubungan
\end{itemize}

\textbf{Strategi untuk Island:}
\begin{enumerate}
    \item \textbf{Reframe Bubble:} Ini bukan kurungan, ini basis operasi untuk menjelajahi dunia
    \item \textbf{Schedule alone time:} Negosiasi waktu sendiri yang jelas dalam Bubble
    \item \textbf{Gradual intimacy:} Perlahan-lahan buka diri, jangan dipaksa
    \item \textbf{Explicit commitment:} Ucapkan komitmen verbal meski terasa canggung
\end{enumerate}
\end{warningbox}

\subsection{Wave (Anxious) Membangun Bubble}

\begin{warningbox}[Tantangan Wave]
\textbf{Hambatan internal Wave:}
\begin{itemize}
    \item Takut Bubble akan pecah kapan saja
    \item Tendensi ``test'' dengan menciptakan konflik
    \item Klaim terlalu banyak waktu pasangan
    \item Sulit percaya pada komitmen pasangan
\end{itemize}

\textbf{Strategi untuk Wave:}
\begin{enumerate}
    \item \textbf{Self-soothing:} Latih menenangkan diri sebelum ``membanjiri'' pasangan
    \item \textbf{Trust the process:} Beri waktu untuk bukti komitmen
    \item \textbf{Communicate needs clearly:} Minta reassurance langsung, tidak lewat test
    \item \textbf{Celebrate security:} Apresiasi saat Bubble terasa kuat
\end{enumerate}
\end{warningbox}

\begin{center}
\renewcommand{\arraystretch}{1.6}
\begin{tabular}{|>{\raggedright}p{2.5cm}|>{\raggedright}p{4cm}|>{\raggedright}p{4cm}|>{\raggedright\arraybackslash}p{4cm}|}
\hline
\rowcolor{primary!20}
\textbf{Tipe} & \textbf{Kontribusi ke Bubble} & \textbf{Yang Perlu Dikerjakan} & \textbf{Partner Support} \\
\hline
\textbf{Anchor} & Stabilitas, kemudahan intimacy, komunikasi jelas & Jangan mengabaikan maintenance & Terima stabilitas, jangan buat drama \\
\hline
\textbf{Island} & Loyalitas, tidak membuat drama, independensi & Buka diri, minta bantuan, prioritaskan hubungan & Beri ruang, jangan chase, tunjukkan komitmen konsisten \\
\hline
\textbf{Wave} & Kedalaman emosional, passion, perhatian & Self-regulate, kurangi testing, percaya komitmen & Beri reassurance, gentle pursuit, validasi emosi \\
\hline
\end{tabular}
\end{center}

\section{Studi Kasus: Strong vs Weak Couple Bubble}

\subsection{Kasus 1: Bubble yang Kuat --- Rina dan Budi}

\begin{insightbox}
\textbf{Profil:} Rina (Wave-ish) dan Budi (Island-ish), menikah 5 tahun, punya 1 anak.

\textbf{Situasi:} Budi baru dipromosikan dengan tanggung jawab besar. Kerjaannya meningkat 60\%. Rina merasa ditinggal dan mulai cemas.

\textbf{Reaksi dengan Bubble Kuat:}
\begin{enumerate}
    \item Rina mengungkapkan perasaannya: ``Aku merasa kehilanganmu akhir-akhir ini. Aku butuh waktu bersamamu.''
    \item Budi \textit{tidak} merasa diserang. Dia mendengar dan validate: ``Aku mengerti. Kerjaan ini memang menguras.''
    \item Mereka sepakat untuk membuat ``protected time'' --- jam 7-9 malam adalah waktu keluarga, tanpa kerjaan.
    \item Budi menyetel batasan di kantor: ``Saya ada jam 7-9, setelah itu saya untuk keluarga.''
    \item Rina bekerja pada self-soothing-nya, tidak membanjiri Budi dengan keluhan.
\end{enumerate}

\textbf{Hasil:} Bubble tetap utuh. Mereka saling melindungi dari tekanan kerja. Rina merasa aman, Budi merasa dihargai.
\end{insightbox}

\subsection{Kasus 2: Bubble yang Lemah --- Dina dan Andi}

\begin{warningbox}
\textbf{Profil:} Dina (Wave) dan Andi (Island), pacaran 3 tahun.

\textbf{Situasi:} Andi sering main game sampai larut malam. Dina merasa diabaikan.

\textbf{Reaksi dengan Bubble Lemah:}
\begin{enumerate}
    \item Dina menunggu dan diam-diam marah (resentmen)
    \item Akhirnya Dina meledak: ``Kamu lebih sayang game daripada aku!''
    \item Andi merasa diserang dan menarik diri: ``Kamu selalu mengeluh. Aku butuh ruang.''
    \item Dina makin panik dan mengejar dengan pesan panjang lebar
    \item Andi semakin menjauh dan menghabiskan waktu dengan teman
    \item Dina curhat ke teman-temannya di media sosial
    \item Andi merasa dikhianati karena masalah mereka dibuat publik
\end{enumerate}

\textbf{Hasil:} Bubble pecah. Pola pursue-withdraw memperburuk keadaan. Interferensi teman (dari curahan Dina) semakin merusak. Akhirnya mereka break up.
\end{warningbox}

\subsection{Analisis Perbandingan}

\begin{center}
\renewcommand{\arraystretch}{1.5}
\begin{tabular}{|>{\raggedright}p{3.5cm}|>{\raggedright}p{5cm}|>{\raggedright\arraybackslash}p{5cm}|}
\hline
\rowcolor{sagegreen!20}
\textbf{Aspek} & \textbf{Bubble Kuat (Rina-Budi)} & \textbf{Bubble Lemah (Dina-Andi)} \\
\hline
Komunikasi & Langsung, jujur, tidak menyerang & Indirect, meledak, menyalahkan \\
\hline
Respons terhadap kebutuhan & Negosiasi, kompromi & Menarik diri atau mengejar \\
\hline
Batasan dengan eksternal & Jelas, disepakati bersama & Kabur atau tidak ada \\
\hline
Repair setelah konflik & Cepat, efektif & Tidak terjadi, makin memburuk \\
\hline
Sense of ``team'' & Kuat, menghadapi masalah bersama & Lemah, satu sama lain dianggap musuh \\
\hline
\end{tabular}
\end{center}

\section{Action Steps: Memulai Hari Ini}

\begin{exercisebox}[Program 7 Hari Membangun Couple Bubble]

\textbf{Hari 1: Inisiasi Komitmen}
\begin{itemize}
    \item Duduk berhadapan dengan pasangan
    \item Diskusikan: ``Apa arti Couple Bubble bagi kita?''
    \item Buat dan tanda tangani Couple Agreement sederhana
\end{itemize}

\textbf{Hari 2: Identifikasi Ancaman}
\begin{itemize}
    \item Masing-masing tulis: Apa yang sering mengganggu waktu/wellbeing kita?
    \item Diskusikan ancaman eksternal (kerja, teknologi, orang lain)
    \item Diskusikan ancaman internal (kebiasaan, attachment issues)
\end{itemize}

\textbf{Hari 3: Membuat Ritual}
\begin{itemize}
    \item Ciptakan morning ritual bersama
    \item Ciptakan evening ritual (stress-reducing conversation)
    \item Setujui ``no-phone time'' pertama kalian
\end{itemize}

\textbf{Hari 4: Praktik Proteksi}
\begin{itemize}
    \item Saat salah satu stres, praktikkan co-regulation
    \item Pasangan yang ``kuat'' hari ini menjadi penyangga
    \item Validasi dan dukung tanpa memberi solusi dulu
\end{itemize}

\textbf{Hari 5: Batasan dengan Dunia Luar}
\begin{itemize}
    \item Praktikkan menolak satu undangan/permintaan yang mengganggu waktu berdua
    \item Matikan notifikasi ponsel selama quality time
    \item Komunikasikan pada keluarga/teman tentang batasan baru kalian
\end{itemize}

\textbf{Hari 6: Review dan Repair}
\begin{itemize}
    \item Evaluasi minggu ini: Apa yang berhasil? Apa yang sulit?
    \item Jika ada konflik, praktikkan repair protocol
    \item Update Couple Agreement jika diperlukan
\end{itemize}

\textbf{Hari 7: Rayakan}
\begin{itemize}
    \item Rayakan komitmen baru kalian (date night spesial)
    \item Tulis appreciation letter untuk pasangan
    \item Rencanakan maintenance ritual untuk minggu-minggu berikutnya
\end{itemize}
\end{exercisebox}

\section{Kesimpulan: Couple Bubble sebagai Investasi}

\begin{insightbox}
\textbf{Key Takeaway:} Couple Bubble bukanlah kemewahan --- ini adalah \highlight{kebutuhan fundamental} untuk hubungan yang sehat dan bertahan lama. Investasi waktu dan energi untuk membangun dan menjaga Bubble akan membayar dengan:

\begin{itemize}
    \item \textbf{Keamanan emosional} untuk berkembang sebagai individu
    \item \textbf{Resiliensi} menghadapi stres dan tantangan hidup
    \item \textbf{Kedalaman koneksi} yang tidak bisa diganggu gugat
    \item \textbf{Kebahagiaan} yang berkelanjutan, bukan sekadar euforia sesaat
\end{itemize}

\textit{Ingatlah: Bubble yang kuat tidak terbentuk dalam semalam. Dibutuhkan komitmen harian, repair saat bocor, dan kesediaan untuk menjaga hubungan sebagai prioritas utama. Tapi hasilnya --- memiliki orang yang selalu ``ada untukmu'' --- sepadan dengan setiap usaha.}
\end{insightbox}

\begin{tipbox}[Reminder Card: Couple Bubble]
\textbf{Tiga Prinsip Utama:}
\begin{enumerate}
    \item \textbf{Prioritas:} Hubungan di atas segalanya
    \item \textbf{Proteksi:} Jadi go-to person satu sama lain
    \item \textbf{Pertahanan:} Lindungi dari ancaman internal dan eksternal
\end{enumerate}

\textbf{Maintenance Harian:}
\begin{itemize}
    \item[$\checkmark$] Morning ritual dengan sentuhan
    \item[$\checkmark$] 20 menit stress-reducing conversation
    \item[$\checkmark$] No-phone zone
    \item[$\checkmark$] Appreciation dan affection
\end{itemize}

\textbf{Saat Ada Masalah:}
\begin{itemize}
    \item Stop pola destruktif
    \item Assess dan komunikasikan
    \item Repair secepatnya
    \item Rebuild dengan ritual
\end{itemize}
\end{tipbox}
